\chapter{MANUAL DE INSTALACIÓN Y DESPLIEGUE}
\label{anx:instalacion}

\section{Requisitos}

\begin{itemize}
    \item Node.js 22 y npm.
    \item Android Studio con el SDK de Android y un emulador, o un dispositivo
    Android con la depuración USB activada.
    \item Una cuenta de Firebase con un proyecto que tenga habilitados
    Authentication (correo y Google), Cloud Firestore y Cloud Storage.
    \item Git, para obtener el código fuente.
\end{itemize}

\section{Obtención del código e instalación de dependencias}

\begin{lstlisting}[language={}, numbers=none, caption={Instalación del proyecto}, label={lst:instalacion}]
git clone https://github.com/ErikaCA99/CoachAppProject.git coach-app
cd coach-app
git checkout develop
npm install
\end{lstlisting}

\section{Configuración de Firebase}

\begin{enumerate}
    \item En la consola de Firebase, registrar una aplicación Android con el
    identificador de paquete definido en \archivo{app.json} y descargar el
    archivo \archivo{google-services.json} en la raíz del proyecto.
    \item Crear el archivo \archivo{.env} en la raíz con las claves de la
    aplicación web del proyecto de Firebase (Código~\ref{lst:env}).
    \item Publicar las reglas de seguridad copiando el contenido de
    \archivo{firestore.rules} y \archivo{storage.rules} en las secciones de
    reglas de Firestore y Storage de la consola (Anexo~\ref{anx:reglas}).
    \item Cargar el catálogo de máquinas ejecutando el \emph{script} de carga
    con las credenciales de una cuenta de servicio, que nunca deben incluirse
    en la aplicación.
\end{enumerate}

\begin{lstlisting}[language={}, numbers=none, caption={Variables de entorno (valores de ejemplo)}, label={lst:env}]
EXPO_PUBLIC_FIREBASE_API_KEY=...
EXPO_PUBLIC_FIREBASE_AUTH_DOMAIN=...
EXPO_PUBLIC_FIREBASE_PROJECT_ID=...
EXPO_PUBLIC_FIREBASE_STORAGE_BUCKET=...
EXPO_PUBLIC_FIREBASE_MESSAGING_SENDER_ID=...
EXPO_PUBLIC_FIREBASE_APP_ID=...
EXPO_PUBLIC_FIREBASE_MEASUREMENT_ID=...
EXPO_PUBLIC_GOOGLE_WEB_CLIENT_ID=...
\end{lstlisting}

\section{Ejecución en desarrollo}

\begin{lstlisting}[language={}, numbers=none, caption={Compilación y ejecución en Android}, label={lst:ejecucion}]
npx expo run:android      # compila e instala el development build
npx expo start            # servidor de desarrollo (en ejecuciones posteriores)
\end{lstlisting}

\section{Verificación del código}

\begin{lstlisting}[language={}, numbers=none, caption={Verificaciones estáticas}, label={lst:verificacion}]
npx tsc --noEmit
npx eslint .
npx prettier --check "src/**/*.{ts,tsx}"
npx expo export -p web
\end{lstlisting}

\section{Generación del instalador}

Con la herramienta EAS de Expo se generan los paquetes de distribución a
partir de los perfiles de \archivo{eas.json}:

\begin{lstlisting}[language={}, numbers=none, caption={Compilación con EAS}, label={lst:eas}]
npm install -g eas-cli
eas login
eas build --platform android --profile preview      # APK de prueba
eas build --platform android --profile production   # AAB para Google Play
\end{lstlisting}
